\subsection{Neural score identities and the limits of conditional inertness}
\label{app:neural-binary-inertness}

Expected binary reward equals the total correct mass; conditional diversity
depends on its allocation among solution modes.
\citet[Theorem~3.1, v1]{sinha2026expected} derive the expected-return
log-ratio dynamics for softmax outcome logits, and
\citet[Section~4.1, v1]{lochab2026ucpo} identify binary reward's
indifference to allocation within the correct set. For differentiable neural
policies, outcome-binary group estimators have per-prompt means proportional
to the correctness gradient but can differ in covariance. Finite steps,
shared prompts, and response-dependent weighting limit the resulting
trajectory equivalence.

\begin{theorem}[Per-prompt neural mean and covariance]
\enforcehalffinalproseline
\label{thm:neural-binary-moments}
Fix a prompt and a differentiable policy $\pi_\theta$ on a finite or countable
response space with common positive support, where $\theta$ lies in an open
subset of $\mathbb R^d$. Let $\mathcal C$ be a fixed verifier-positive event,
$P=\Pr_\theta(Y\in\mathcal C)\in(0,1)$, and
$s(Y)=\nabla_\theta\log\pi_\theta(Y)$. Assume differentiation may be
interchanged with the probability sums over the full support and over
$\mathcal C$, and that $\E\|s(Y)\|_2^2<\infty$. For $G$ independent
responses, write $R_i=\1[Y_i\in\mathcal C]$, $R=\sum_iR_i$, and let
$a:\{0,1\}\times\{0,\ldots,G\}\to\mathbb R$ be any finite, detached
advantage. Define
\[
 \widehat g_a=\frac1G\sum_i a(R_i,R)s(Y_i),\qquad
 v=\nabla_\theta P,\qquad
 \Sigma_r=\operatorname{Cov}\!\left(s(Y)\mid\1[Y\in\mathcal C]=r\right).
\]
With
\[
 \beta_r=\frac1G\left(\frac{r\,a(1,r)}P
          -\frac{(G-r)a(0,r)}{1-P}\right),
 \qquad c_a(P)=\E[\beta_R],
\]
where $R\sim\operatorname{Binomial}(G,P)$, the exact identities are
\begin{align}
 \E[\widehat g_a]&=c_a(P)v,
 \label{eq:neural-binary-mean}\\
 \operatorname{Cov}(\widehat g_a)
 &=\operatorname{Var}(\beta_R)vv^\top
   +\frac1{G^2}\E\!\left[
       R\,a(1,R)^2\Sigma_1+(G-R)a(0,R)^2\Sigma_0\right].
 \label{eq:neural-binary-covariance}
\end{align}
In particular, $c_a$ is the coefficient in
Equation~\eqref{eq:objective-inertness}, without an independent-logit
assumption.
\end{theorem}

\begin{proof}
\enforcehalffinalproseline
The conditional-score identity
\citep[Theorem~1]{tajwar2026maxrl} (also
\citealp[Lemma~B.1]{avspo2026}) applied to the correct and incorrect events gives
\[
 \E[s(Y)\mid Y\in\mathcal C]=v/P,
 \qquad
 \E[s(Y)\mid Y\notin\mathcal C]=-v/(1-P).
\]
Conditioned on the complete reward vector, the response scores remain
independent. The conditional mean of $\widehat g_a$ is $\beta_Rv$ and
its conditional covariance is
$G^{-2}[R a(1,R)^2\Sigma_1+(G-R)a(0,R)^2\Sigma_0]$.
Taking expectations proves Equation~\eqref{eq:neural-binary-mean}; the law
of total covariance proves Equation~\eqref{eq:neural-binary-covariance}.
\end{proof}

For one prompt, a common locally Lipschitz preconditioner $M(\theta)$ and
positive coefficients $c_a(P)$ give deterministic mean fields
$\dot\theta=c_a(P)M(\theta)\nabla_\theta P$ that trace the same orbit
under a positive change of time, wherever the flows exist uniquely.
This is a statement about the mean direction. Equation~\eqref{eq:neural-binary-covariance}
shows that outcome-binary estimators can have different sampling covariance,
including after their mean magnitudes are matched.

\begin{corollary}[A finite-step comparison in neural parameters]
\enforcehalffinalproseline
\label{cor:neural-binary-finite-step}
Under Theorem~\ref{thm:neural-binary-moments}, suppose $c_a(P)>0$ and take
$\theta^+=\theta+\eta\widehat g_a/c_a(P)$ with $\eta>0$. Let $f$ be a
twice differentiable scalar observable with
$\|\nabla^2f\|_{\mathrm{op}}\le L_f$ along every update segment, which
is assumed to remain in its domain. Then, conditional on the current policy,
\begin{equation}
 \left|\E[f(\theta^+)-f(\theta)]
       -\eta\nabla f(\theta)^\top v\right|
 \le\frac{L_f\eta^2}{2}
 \left(\|v\|_2^2+
       \frac{\operatorname{tr}\operatorname{Cov}(\widehat g_a)}{c_a(P)^2}
 \right).
 \label{eq:neural-binary-finite-step}
\end{equation}
\end{corollary}

\begin{proof}
\enforcehalffinalproseline
Taylor's theorem bounds the remainder after the linear term by
$L_f\eta^2\|\widehat g_a\|_2^2/(2c_a(P)^2)$. Take expectations and use
Equation~\eqref{eq:neural-binary-mean} together with
$\E\|\widehat g_a\|_2^2
=c_a(P)^2\|v\|_2^2+\operatorname{tr}\operatorname{Cov}(\widehat g_a)$.
\end{proof}

For two estimators satisfying the corollary, division by their respective
coefficients $c_a(P)$ gives the same mean step $\eta v$; the training updates
do not use this normalization. Their expected changes in $f$ then differ by
at most the sum of the two second-order bounds. Conditional solution mode
probabilities or \pmd{} may be used as $f$ where the stated smoothness bounds
hold. Mean-direction equivalence is
therefore a first-order statement for stochastic finite updates; it does
not determine their trajectories or eventual surviving solution modes.
The bound assumes fresh on-policy score updates and does not cover Adam's
moment estimates or subsequent PPO epochs.

\begin{remark}[Shared prompts need not follow the same trajectory]
\label{rem:neural-binary-shared-prompts}
For a fixed prompt-sampling distribution $\mu_x$ and shared parameters,
the aggregate mean, assuming integrable updates, is
\begin{equation}
 \E[\widehat g_a]
 =\sum_x\mu_x c_a(P_x)\nabla_\theta P_x.
 \label{eq:neural-binary-prompt-sum}
\end{equation}
Different objectives generally assign different relative coefficients to
the prompt gradients, so the per-prompt identity does not imply a common
aggregate orbit. An explicit example also shows a difference in conditional
diversity. Let $\theta=(u,v)$, $\sigma(t)=(1+e^{-t})^{-1}$, and sample two
prompts $A,B$ equally. Set
\[
 P_A=\sigma(u+v),\qquad P_B=\sigma(u-v),\qquad
 q_A=(\sigma(v),1-\sigma(v)).
\]
For prompt $A$, the probabilities of two correct responses and one incorrect
response are $(P_A\sigma(v),P_A[1-\sigma(v)],1-P_A)$; for $B$, use
$(P_B/2,P_B/2,1-P_B)$. At $(u,v)=(0,\log3)$,
$P_A=3/4$, $P_B=1/4$, and
$\nabla P_A=(3/16,3/16)$, $\nabla P_B=(3/16,-3/16)$.
For $G=3$, Lemmas~\ref{lem:grpo-mean} and~\ref{lem:maxrl-mean} give
$c_{\mathrm{Dr}}=2/3$ and $c_{\mathrm{Max}}(P)=2-P$, hence
\[
 \dot\theta_{\mathrm{Dr}}=(1/8,0),\qquad
 \dot\theta_{\mathrm{Max}}=(9/32,-3/64).
\]
Prompt $A$ has conditional disagreement
$D_A=1-\sum_cq_{A,c}^2=2\sigma(v)[1-\sigma(v)]$, so direct
differentiation gives
\[
 \dot D_{A,\mathrm{Dr}}=0,\qquad
 \dot D_{A,\mathrm{Max}}=9/1024>0.
\]
Both prompts' correctness increases under both mean fields, since
$\dot P_A=(3/16)(\dot u+\dot v)>0$ and
$\dot P_B=(3/16)(\dot u-\dot v)>0$ at this point. Changing a binary
objective can therefore change conditional diversity through shared prompt
gradients, despite the exact per-prompt identity.
\end{remark}

\begin{proposition}[The residual from response-dependent normalization]
\enforcehalffinalproseline
\label{prop:neural-binary-length-residual}
Under Theorem~\ref{thm:neural-binary-moments}, multiply each response's score
by a detached function $\omega(Y)$ with finite conditional second moments;
inverse response length is the case $\omega(Y)=1/L(Y)$ with $L(Y)>0$. Define
\[
 \widehat g_{a,\omega}
 =\frac1G\sum_i a(R_i,R)\omega(Y_i)s(Y_i),\qquad
 A_1=\frac{\E[R a(1,R)]}{G},\quad
 A_0=\frac{\E[(G-R)a(0,R)]}{G},
\]
and, conditioning on $\1[Y\in\mathcal C]=r$, let
$\bar\omega_r=\E[\omega(Y)\mid r]$,
$K_r=\operatorname{Cov}(\omega(Y),s(Y)\mid r)$, and
$\sigma_{\omega,r}^2=\operatorname{Var}(\omega(Y)\mid r)$.
Then
\begin{align}
 \E[\widehat g_{a,\omega}]
 &=\widetilde c_a\nabla_\theta P+A_1K_1+A_0K_0,
 &\widetilde c_a&=\frac{A_1\bar\omega_1}{P}
                  -\frac{A_0\bar\omega_0}{1-P},
 \label{eq:neural-binary-length-residual}\\
 \|A_1K_1+A_0K_0\|_2
 &\le\sum_{r=0}^1|A_r|\sigma_{\omega,r}
                 \sqrt{\operatorname{tr}\Sigma_r}.
 \label{eq:neural-binary-length-bound}
\end{align}
\end{proposition}

\begin{proof}
\enforcehalffinalproseline
Conditioning on the reward vector gives
$\E[\widehat g_{a,\omega}]
=A_1\E[\omega s\mid1]+A_0\E[\omega s\mid0]$.
For each outcome, $\E[\omega s\mid r]
=\bar\omega_r\E[s\mid r]+K_r$; substituting the conditional score
means proves Equation~\eqref{eq:neural-binary-length-residual}.
The vector Cauchy--Schwarz inequality gives
$\|K_r\|_2\le\sigma_{\omega,r}\sqrt{\operatorname{tr}\Sigma_r}$;
the triangle inequality completes the bound.
\end{proof}

The residual is a within-outcome covariance between response weighting and
the policy score. For inverse-length weighting, a common response length
makes it zero; equal mean lengths for correct and incorrect responses need
not. Response-dependent weights need not be functions of binary outcomes
alone. IPS uses outcome-frequency weights
\citep[Section~4.4, v1]{sinha2026expected}, and UCPO adds a
conditional-uniformity term \citep[Section~6, v1]{lochab2026ucpo}; these
interventions are outside the outcome-binary class analyzed above.

The verified-key diversity metric is a specialization of SetPO's
distributional functional. Its influence gives an exact metric gradient
for shared neural parameters and fixed verified execution keys.

\begin{corollary}[Verified-key specialization of SetPO]
\enforcehalffinalproseline
\label{cor:pcmd-setpo-neural}
Fix a prompt, a verifier-positive response event $\mathcal C$, and a
deterministic key map $V:\mathcal C\to\{1,\ldots,m\}$, all independent
of $\theta$. Let $\pi_\theta$ be differentiable with common positive support,
$P_\theta=\Pr_\theta(\mathcal C)>0$, and
$Q_\theta=\pi_\theta(\cdot\mid\mathcal C)$. Assume the score
$s_\theta(Y)=\nabla_\theta\log\pi_\theta(Y)$ is integrable under $Q_\theta$
and differentiation commutes with the response sums over each key.
Write $q_c=Q_\theta(V(Y)=c)$ and $S_2=\sum_cq_c^2$.
For $k(y,y')=\1[V(y)=V(y')]$ and $g(u)=1-u$, the functional
$\mathcal F(Q)=\E_{Y\sim Q}[g(\E_{Y'\sim Q}k(Y,Y'))]$ is exactly
$\pmd(q)$. The influence formula of
\citet[Theorem~4.2, v1]{li2026setpo} specializes, for $y\in\mathcal C$, to
\begin{align}
 \left.\frac{d}{d\varepsilon}
  \mathcal F((1-\varepsilon)Q_\theta+\varepsilon\delta_y)
  \right|_{\varepsilon=0}
 &=2(S_2-q_{V(y)}),\label{eq:pcmd-setpo-influence}\\
 \nabla_\theta\pmd(q)
 &=2\E_{Y\sim Q_\theta}
       [(S_2-q_{V(Y)})s_\theta(Y)].
 \label{eq:pcmd-neural-gradient}
\end{align}
\end{corollary}
\begin{proof}
\enforcehalffinalproseline
The equality kernel and $g$ satisfy SetPO's Assumption~4.1, and its local
mass is $q_{V(y)}$. Substitution into its influence formula proves the first
identity. The influence has mean zero under $Q_\theta$, so the
$-\nabla_\theta\log P_\theta$ term in the conditional score cancels,
giving the second identity.
\end{proof}
These identities permit shared neural parameters. Along an update, the
sign of the diversity change depends on its alignment with this metric
gradient; the influence formula alone supplies no monotonicity guarantee
for a given training direction.

